% Information Security & Acceptable Use Policy.
% Content from "Information Security & Acceptable Use Policy.docx".
\documentclass{paridyum}
\begin{document}

\paridyumletterhead
\paridyumdateref{\dkDocDate}{\dkDocRef}
\paridyumtitle{Information Security \& Acceptable Use Policy}

\textit{This policy applies to interns who are given access to Paridyum's accounts,
source code, documents, systems, devices, infrastructure, or other digital resources.}

\section*{1. Purpose}
The purpose of this policy is to protect Paridyum's systems, source code, confidential
information, credentials, and other digital assets. Interns must use Paridyum resources
responsibly and only for authorized internship activities.

\section*{2. Accounts and Passwords}
The Intern must:
\begin{itemize}[leftmargin=*, noitemsep]
    \item keep passwords and credentials private;
    \item use strong and unique passwords;
    \item enable multi-factor authentication where required;
    \item never share Company accounts with another person;
    \item never use another person's account without authorization; and
    \item immediately report lost, stolen, or compromised credentials.
\end{itemize}

\section*{3. Source Code and Repositories}
Paridyum's private source code must remain within authorized repositories and systems.
The Intern must not:
\begin{itemize}[leftmargin=*, noitemsep]
    \item publish private repositories;
    \item upload Company code to personal public repositories;
    \item share source code with unauthorized persons;
    \item copy Company code for unrelated personal projects; or
    \item intentionally remove or bypass repository security controls.
\end{itemize}

\section*{4. API Keys and Secrets}
Passwords, API keys, access tokens, private keys, database credentials, and other
secrets must never be:
\begin{itemize}[leftmargin=*, noitemsep]
    \item committed directly into source code;
    \item uploaded to public repositories;
    \item shared through unauthorized services;
    \item posted in public communication channels; or
    \item shared with another person without authorization.
\end{itemize}
If a secret is accidentally exposed, the Intern must report it immediately.

\section*{5. Access to Systems}
The Intern may only access systems, repositories, databases, files, and information
required for assigned internship responsibilities. The Intern must not attempt to:
\begin{itemize}[leftmargin=*, noitemsep]
    \item access restricted systems;
    \item bypass security controls;
    \item obtain unauthorized administrator privileges;
    \item access another user's account;
    \item view unrelated confidential information; or
    \item perform unauthorized security testing.
\end{itemize}

\section*{6. Personal Devices}
If a personal laptop or device is used for internship work, the Intern should:
\begin{itemize}[leftmargin=*, noitemsep]
    \item use a password, PIN, or biometric device lock;
    \item keep the operating system reasonably updated;
    \item take reasonable precautions against malware;
    \item prevent unauthorized persons from accessing Company information; and
    \item securely remove Paridyum information when requested at the end of the
          internship.
\end{itemize}

\section*{7. AI Tools}
The Intern must not provide confidential Paridyum information to external AI tools
unless their use has been approved. Do not submit the following to unauthorized AI
services:
\begin{itemize}[leftmargin=*, noitemsep]
    \item private source code;
    \item passwords or credentials;
    \item API keys;
    \item customer or user information;
    \item confidential documents;
    \item unreleased product information;
    \item internal business information; or
    \item sensitive technical architecture.
\end{itemize}
AI-generated code or content must be reviewed before being added to Paridyum projects.

\section*{8. Third-Party Services}
The Intern must not upload Paridyum information or source code to unauthorized:
\begin{itemize}[leftmargin=*, noitemsep]
    \item cloud-storage services;
    \item code-sharing websites;
    \item file-sharing services;
    \item online development environments;
    \item collaboration platforms; or
    \item other third-party services.
\end{itemize}
If a new service is required for the internship, approval should be obtained before
using it with Paridyum information.

\section*{9. Software Installation}
Software used for Paridyum work should come from legitimate and reasonably trusted
sources. The Intern must not knowingly install or use:
\begin{itemize}[leftmargin=*, noitemsep]
    \item pirated software;
    \item malware;
    \item credential-stealing tools;
    \item unauthorized remote-access software; or
    \item software intended to bypass Paridyum's security controls.
\end{itemize}

\section*{10. Security Testing}
The Intern must not perform penetration testing, vulnerability scanning, exploitation,
load testing, or similar security testing against Paridyum systems unless specifically
authorized. Any authorized security testing must remain within the approved scope.

\section*{11. Confidential Information}
Confidential information must be handled according to the \textbf{Paridyum
Non-Disclosure Agreement}. The Intern must take reasonable precautions to prevent
confidential information from being accidentally viewed, copied, published, or shared.

\section*{12. Security Incidents}
The Intern must immediately report suspected security incidents, including:
\begin{itemize}[leftmargin=*, noitemsep]
    \item leaked passwords or API keys;
    \item accidental public repositories;
    \item suspicious account activity;
    \item unauthorized access;
    \item malware;
    \item lost or stolen devices containing Paridyum information;
    \item accidental disclosure of confidential information; or
    \item discovered security vulnerabilities.
\end{itemize}

\vspace{0.3em}
\begin{infobox}
\textbf{Report security issues to}\\[0.4em]
\begin{tabularx}{\linewidth}{@{}>{\bfseries}p{2.2cm}X@{}}
Contact: & \dkOrgSecurityContact{} --- \dkOrgSecurityContactRole \\
Email:   & \dkOrgSecurityEmail
\end{tabularx}
\end{infobox}

\vspace{0.3em}
The Intern should report accidental incidents promptly rather than attempting to hide
or independently resolve a serious security issue.

\section*{13. Company Equipment}
Any equipment provided by Paridyum remains the property of Paridyum unless otherwise
agreed in writing. Company equipment must:
\begin{itemize}[leftmargin=*, noitemsep]
    \item be used responsibly;
    \item be reasonably protected from loss or damage;
    \item not be given to another person without permission; and
    \item be returned when requested or when the internship ends.
\end{itemize}

\section*{14. End of Internship}
When the internship ends, the Intern must:
\begin{itemize}[leftmargin=*, noitemsep]
    \item stop accessing Paridyum systems;
    \item return Company equipment;
    \item remove Company credentials from personal devices;
    \item return or delete confidential files when requested; and
    \item cooperate with account and access removal.
\end{itemize}
The Intern must not retain Paridyum's private source code, credentials, confidential
documents, or other protected information after the internship.

\section*{15. Violations}
Serious or repeated violations of this policy may result in removal of system access or
termination of the internship. Where appropriate, Paridyum may also take other action
available under applicable law.

\section*{16. Acknowledgement}
By signing below, the Intern confirms that they have read and understood this
Information Security \& Acceptable Use Policy and agree to follow it during their
internship.

\vspace{1.5em}
\paridyumsignblock

\vspace{1.5em}
\paridyumcountersign{Intern}{\dkPIdentityLegalName}

\end{document}
